\section{The physical model and prescribed initial data}
\label{sec:model}

\subsection{An independent belt parameter}
The oscillator consists of two point masses joined by a rigid link.
The upper mass moves vertically against a spring and damper; the lower
mass can contact a belt moving to the left. In the nondimensional
coordinates of \cite[Section~6.1]{NVC2018}, both masses and the link
length are one. The link angle $\phi$ is measured from the downward
vertical, and the upper mass is at height $1+\psi$ above the undeformed
belt. A torsional spring acts on $\phi$. The observables $y$ and $u$
below are the lower mass's height and its slip relative to the belt.
Let $q=(\phi,\psi)^T$, $\omega=\dot\phi$, and define
\begin{equation}
 y=\psi+1-\cos\phi,\qquad v=\dot y,
 \qquad u=\omega\cos\phi+U_b,\qquad U_b=3/5+\eta.
 \label{eq:bd-observables}
\end{equation}
The belt detuning $\eta$ is independent of the mechanical parameters.
We use the critical mechanical family of \cite[Section~6]{NVC2018}:
\begin{gather}
 M(\phi)\ddot q=f_{\beta,\kappa}(q,\dot q)
                  +Q_N(\phi)\lambda_N+Q_T(\phi)\lambda_T,
 \label{eq:bd-mechanics}\\
 M=\begin{pmatrix}1&\sin\phi\\\sin\phi&2\end{pmatrix},
 \quad Q_N=\binom{\sin\phi}{1},\quad Q_T=\binom{\cos\phi}{0},
 \nonumber\\
 f_{\beta,\kappa}=
 \binom{-(\phi-\phi_0)-\sin\phi}
       {-\kappa\psi-c_\psi\dot\psi-2-\cos\phi\,\omega^2},
 \nonumber\\
 \mu=25/12,\quad c_\psi=25(18-7\beta)/108,
 \quad\phi_0=\phi^*+49/20-7\beta/12-9\kappa/50,
 \nonumber\\
 \quad (\cos\phi^*,\sin\phi^*)=(3/5,4/5).
 \label{eq:bd-critical-family}
\end{gather}
The prescribed completion is
\begin{equation}
 \eps\dot z=-z-\eps^2\lambda_N,\qquad
 \lambda_N=\eps^{-2}(z-y)_+,\qquad
 \lambda_T=-\mu\Phi_\theta(u/\eps)\lambda_N.
 \label{eq:bd-completion}
\end{equation}
Thus contact is determined by $g=y-z$, not by $y$ alone.
In flight $z(t)=z_r e^{-(t-t_r)/\eps}$, and
$\dot g=v+z/\eps$ at release. Every release below is the first
zero of the active normal force on the stated continuous contact
passage, and every touchdown is the first zero of the subsequent
positive free gap. Initial capture may contain other contacts.

Write $X=(\phi,y,\omega,v)$ and eliminate
$\psi=y-1+\cos\phi$, $\dot\psi=v-\sin\phi\,\omega$.
The free mechanical vector field and saturated contact field per unit
normal force are
\begin{equation}
 F_{\rm free}=(\omega,v,(M^{-1}f)_1,b),\qquad
 Q_{\rm sat}=(0,0,(M^{-1}(Q_N-\mu Q_T))_1,p),
 \label{eq:bd-fields}
\end{equation}
where
\begin{equation}
 b=Q_N^TM^{-1}f+\cos\phi\,\omega^2,\qquad
 a_T=Q_T^TM^{-1}f-\sin\phi\,\omega^2,
 \qquad p=C-\mu B.
 \label{eq:bd-free-accelerations}
\end{equation}
The exact contact equations for the observables are
\begin{equation}
 \dot v=b+(C-\mu B\Phi_\theta(u/\eps))\lambda_N,
 \qquad
 \dot u=a_T+(B-\mu A\Phi_\theta(u/\eps))\lambda_N.
 \label{eq:bd-contact-accelerations}
\end{equation}
Direct inversion of $M$ gives
\begin{equation}
 A=\frac{2\cos^2\phi}{1+\cos^2\phi},\qquad
 B=\frac{\sin\phi\cos\phi}{1+\cos^2\phi},\qquad
 C=\frac1{1+\cos^2\phi},\qquad
 D=C-B^2/A=\frac12.
 \label{eq:bd-delassus}
\end{equation}
In particular the coefficient after eliminating tangential force is
positive even when the saturated coefficient $p$ is zero or negative.
The chart is restricted throughout so that $\cos\phi$ has a common
positive lower bound.
Since $p=C(\phi)-\mu B(\phi)$, the rate
$p_\phi(G)=a$ below permits a further common restriction with
$p_\phi>0$. Each observation section $p=-as$ then fixes one angle
$\phi_s$ on this branch.

\subsection{Critical rates and the working box}
At $G=(\phi^*,0,-1,0)$ in the mechanical coordinates,
$p=b=0$ and the tuned slip is zero. The detuned slip at $G$ is
$\eta$. Substitution in \eqref{eq:bd-fields}--\eqref{eq:bd-delassus}
gives
\begin{gather}
 A_G=9/17,\quad B_G=6/17,\quad C_G=25/34,\quad
 h:=\mu A_G-B_G=3/4,\qquad a:=175/408,
 \nonumber\\
 F_{\rm free}p(G)=-a,\quad F_{\rm free}b(G)=-c_\lambda,
 \quad a_T(G)=d_\lambda,\quad Q_{\rm sat}b(G)=a\beta,
 \label{eq:bd-critical-rates}\\
 c_\lambda=\frac{126765+61425\beta-30625\beta^2
                    -(9450\beta+8262)\kappa}{55080},\qquad
 d_\lambda=\frac{495-175\beta-54\kappa}{180},
 \nonumber\\
 \ell:=\lambda_*=\frac{c_\lambda}{a(\beta-1)},\qquad
 g_*:=d_\lambda-h\ell
       =\frac{55825\beta+17712\kappa-213390}{31500(\beta-1)}.
 \label{eq:bd-incoming-rates}
\end{gather}
Here $\lambda=(\beta,\kappa)$; we also write $c=c_\lambda$
and $\gamma=-g_*>0$. These formulas can be checked by differentiating the
explicit fields in \eqref{eq:bd-fields}. At $G$ the added generalized
acceleration is $(-5/4,1)$ in $(\ddot\phi,\ddot\psi)$, so its
$(\dot\omega,\dot v)$ components are $(-5/4,0)$, and
\begin{equation}
 p_\phi(G)=a,\qquad
 \det D_{(\phi,\omega)}(p,b)(G)=-6125\beta/41616.
 \label{eq:bd-chart-determinant}
\end{equation}
Thus $(p,b,y,v)$ are analytic coordinates on a common neighborhood;
write $\Xi(p,b,y,v)$ for their inverse mechanical chart.

Fix a sufficiently small closed box $\Lambda$ about $(7/3,0)$ in
the half-plane $\kappa\ge0$, with
\begin{equation}
 2<\beta_-\le\beta\le\beta_+<5/2,
 \qquad c,d_\lambda,\ell,-g_*>0,
 \qquad |\det D_{(\phi,\omega)}(p,b)(G)|>0,
 \label{eq:bd-box}
\end{equation}
all with common strict margins. Uniform statements use this compact
box; its interior is relative to $\kappa\ge0$.
The upper bound $\beta_+<5/2$ records the scope of the present proof.
The coefficient $\beta$ is also the exponent in the incoming weak
trace, as the scalar equation in Section~\ref{sec:entry}
shows. We use the prescribed compact range \eqref{eq:bd-box} to
construct the weak curve and obtain the common forcing bound
\eqref{ent:forcing-domination} in the upstream normal-energy and
weak-transport estimates. The strict sign and chart margins fix the
event orientation and the common coordinate neighborhood.

\subsection{Profiles and physical preparation}
Let $\Theta$ be compact, and let $\theta\mapsto\Phi_\theta$ be
continuous in the metric
\begin{equation}
 d_{\bar\alpha}(\Phi,\Psi)=
 \max_{0\le j\le4}\sup_{q\in\R}
 e^{\bar\alpha|q|}|\partial_q^j(\Phi-\Psi)(q)|,
 \qquad 0<\alpha<\bar\alpha<2.
 \label{eq:bd-profile-metric}
\end{equation}
Its image is compact in the affine space consisting of $\tanh$ plus
functions with finite weighted norm. Fix a compact
$I_s=[s_-,s_+]\Subset(-1,1)$ containing $B/(\mu A)$ on the
working mechanical chart. The following bounds hold with common
constants $0<M_{\rm sl}<Q_\Phi$, $m_\Phi,\varkappa_{\rm sl},C_\Phi>0$:
\begin{enumerate}[label=\textup{(F\arabic*)}]
\item $\Phi_\theta$ is strictly increasing, $\Phi_\theta(0)=0$,
 has values in $(-1,1)$ and limits $\pm1$ at $\pm\infty$.
\item $\Phi_\theta([-Q_\Phi,Q_\Phi])\supset I_s$,
 $\Phi_\theta'\ge m_\Phi$ there, and its first four derivatives
 are uniformly bounded there.
\item $\Phi_\theta(M_{\rm sl})\ge s_++\varkappa_{\rm sl}$ and
 $\Phi_\theta(-M_{\rm sl})\le s_--\varkappa_{\rm sl}$.
\item For $q\ge M_{\rm sl}$,
 \[
 1-\Phi_\theta(q)+\sum_{j=1}^4q^j|\Phi_\theta^{(j)}(q)|
                 \le C_\Phi e^{-\alpha q},
 \]
 with the reflected bound on the negative tail.
\end{enumerate}
The slip threshold $M_{\rm sl}$ is fixed with the profile class and
mechanical chart.
Nonodd profiles are allowed. Each profile is fixed before taking
the stiff limit. The family $\tanh(\theta q)$, $1\le\theta\le5/4$,
is an example: all derivatives through order four are bounded by
$C(1+|q|^4)e^{-2|q|}$ on the tails, as are the label derivatives
up to an extra polynomial factor. This proves both (F4) and
continuity in \eqref{eq:bd-profile-metric}; monotonicity gives the
remaining conditions after the strip is chosen large enough.

Choose a small fixed $p_0>0$ and a continuous compact family of
saturated rigid data $(\phi_\nu,\omega_\nu)$ on $p=p_0$.
Their positive-slip rigid orbits reach $G$ and have traces
\begin{equation}
 b_{{\rm rig},\nu}(p)=\bar b_\lambda(p)+C_1(\lambda,\nu)p^\beta
                 +\mathcal R_C(p),\qquad
 |\partial_p^j\mathcal R_C(p)|\le Cp^{\beta+1-j},\quad 0\le j\le3,
 \label{eq:bd-prescribed-trace}
\end{equation}
where $\bar b_\lambda(p)=-\ell p+O(p^2)$ is the analytic weak
curve constructed in Section~\ref{sec:entry}.
On each closed rigid subarc before $G$, the force $-b/p$ and tuned
slip have common positive lower bounds. The rigid trace is unchanged
by detuning: $U_b$ enters only the slip observable and the profile.

For a fixed fast triple $\mathbf n=(Y_0,V_0,Z_0)$ in any prescribed
compact set, the complete physical initial state is
\begin{equation}
 \cX_\eps(0;\lambda,\theta,\nu,\mathbf n)
       =(\phi_\nu,\eps^2Y_0,\omega_\nu,\eps V_0,\eps^2Z_0)
 \label{eq:bd-physical-initial-data}
\end{equation}
in coordinates $(\phi,y,\omega,v,z)$.
The slow datum and fast triple stay fixed as $\eps$ varies.
The initial state need not be in contact. For $\mathbf n=0$ it is
the same full physical state for every compliance and profile.
No later release or touchdown state is independently prepared.
The class is nonempty: the analytic weak curve supplies a rigid
datum at $p_0$, and a compact interval of nearby weak coefficients
supplies more. The scalar construction in Section~\ref{sec:entry} gives
\eqref{eq:bd-prescribed-trace}. The positive limits of $-b/p$ and
$u_0/p$ then give the closed-arc margins after the chart is reduced.

\paragraph{Standing convention.}
Fix the compact profile and physical-preparation classes above and
$k\in I=[k_-,k_+]\Subset(0,\infty)$; the nominal mechanical box
may then be reduced jointly for these same classes. All error
estimates and event margins are uniform on the resulting compact
label classes. Unless explicitly stated otherwise, mechanical
parameters and each sufficiently small positive observation section
are fixed before $\eps\downarrow0$. The common graph, return and
section cutoffs are chosen once before these targets. Uniform
selection is asserted only where stated; each incoming datum keeps
its own prescribed rigid arc.

The complete label is $\mathfrak l=(\lambda,k,\theta,\nu,\mathbf n)$.
We use $\cX=(\phi,y,\omega,v,z)$ for the full physical state and
$X=(\phi,y,\omega,v)$ for its mechanical part. In a five-state norm,
$G$ includes its zero memory coordinate. Put
$r=k\sqrt\eps=-\eta$. The actual belt slip is
$u_r=u=\omega\cos\phi+U_*-r$; the critical-speed slip is
$u_0=\omega\cos\phi+U_*=u_r+r$.
The coefficient $k$ and $r$ remain constant along each trajectory;
the physical belt speed follows $U_b=U_*-k\sqrt\eps$ as compliance
varies. The slow datum and scaled fast triple $\mathbf n$ remain fixed;
the physical initial state follows \eqref{eq:bd-physical-initial-data}.
